\section{Justificación}

\subsection{Importancia y Necesidad del Estudio}

La justificación se basa en tres factores que constituyen un solo argumento. Primero,
la brecha en la cobertura de monitoreo emocional en la industria es crítica y
cuantificable: solo el 31\,\% de los contact centers monitorea activamente las
emociones del cliente \cite{zoom2025}, lo que evidencia una necesidad concreta de
herramientas automáticas que amplíen sin incrementar los costos operativos. Segundo,
la tecnología SER ha alcanzado suficiente madurez en condiciones reales: precisiones
del 75--85\,\% en habla espontánea \cite{abdelhamid2022, wagner2023}, superando la
línea de viabilidad para aplicaciones industriales. Tercero, el costo humano del
problema, con tasas de burnout superiores al 63\,\% \cite{sqm2023} y rotaciones
anuales entre el 38\,\% y el 52\,\% \cite{hiringbranch2024}, justifica la inversión en
soluciones preventivas desde una perspectiva ética y económica.

El estudio es, además, necesario porque aborda una limitación de los enfoques
existentes: la dependencia del análisis de texto. Los sistemas de análisis de
sentimiento textual, ampliamente desplegados en la industria, son insensibles a las
señales paralingüísticas que revelan estados emocionales encubiertos. Jha et al.
(2022) demuestran que la combinación de características prosódicas y espectrales
supera consistentemente a los enfoques puramente léxicos en escenarios de habla
natural \cite{jha2022}.

\subsection{Beneficios Esperados e Impacto}

En el ámbito \textbf{operacional y económico}, la automatización de la detección de
picos emocionales permitirá a los supervisores focalizar su atención en el subconjunto
de llamadas con escalamiento real, lo que multiplicará la cobertura efectiva de
monitoreo sin aumentar el personal. Si el sistema alcanza las métricas propuestas
—precisión $\geq 75$\,\%, recall $\geq 70$\,\%— podrá identificar automáticamente los
momentos críticos en el 100\,\% de las llamadas procesadas, en contraste con el
1--5\,\% que hoy es posible revisar manualmente. Esto representa una mejora de al
menos $\times 20$ en la cobertura de la evaluación de la calidad.

En el ámbito del \textbf{bienestar laboral}, la detección temprana de picos
emocionales habilitará intervenciones supervisoras oportunas que pueden interrumpir la
escalada antes de que el agente alcance el umbral de agotamiento emocional. La
literatura señala que el acceso a la retroalimentación inmediata y al apoyo del
supervisor durante interacciones difíciles es uno de los factores protectores más
eficaces contra el burnout \cite{sqm2023}.

En el ámbito \textbf{académico y tecnológico}, el proyecto genera contribuciones en
tres frentes: (a) comparación empírica de tres métodos de cálculo de línea base
emocional individual sobre llamadas reales en español, dominio subrepresentado en la
literatura \cite{jha2022, delope2023}; (b) evaluación comparativa de algoritmos de
detección de anomalías (Z-scores, Isolation Forest, One-Class SVM) aplicados a
features paralingüísticas en conversaciones espontáneas de call center; y (c) un
pipeline reproducible y documentado que puede ser adoptado y extendido por otros
investigadores o desplegado directamente en entornos industriales.

\subsection{Pertinencia con el Conocimiento Actual y los Objetivos}

Esta investigación se inserta en una tendencia clara en el campo del procesamiento de
habla: el desplazamiento desde modelos entrenados en datos actuados hacia sistemas
capaces de operar sobre habla espontánea en dominios reales. Feng y Devillers (2023)
identifican precisamente los call centers de atención al cliente como uno de los
entornos más demandantes y menos atendidos por la literatura de SER en condiciones
reales \cite{feng2023}. Wagner et al. (2023) señalan el cierre de la brecha de
valencia como el desafío central del campo \cite{wagner2023} y De Lope y Graña (2023)
confirman que la generalización a dominios espontáneos sigue siendo el reto abierto
más relevante \cite{delope2023}.

La justificación se articula directamente con los cinco objetivos específicos. El
Objetivo~1 (extracción de características acústicas) responde a la necesidad de
capturar las señales que el análisis de texto omite. El Objetivo~2 (evaluación de
métodos de línea base) aborda la causa metodológica del problema. El Objetivo~3
(detección de anomalías) provee el mecanismo central de identificación de picos. El
Objetivo~4 (clasificación de valencia) añade información sobre la naturaleza del
pico, habilitando intervenciones diferenciadas. El Objetivo~5 (validación con llamadas
reales) asegura que los resultados sean relevantes en el dominio específico de interés,
no solo en benchmarks de laboratorio.
